% I. Предметна област, стандарти, заплахи и избор на решения.
% Обединява предишните раздели „литературен преглед“ и „алтернативи“.
\section{Предметна област и избор на решения}
\label{sec:domain}

\subsection{Стандарти, заплахи и протоколи}
\label{subsec:standards}

Рамката за управление на информационната сигурност идва от серията ISO/IEC 27000. Пряко
приложима тук е ISO/IEC 27033-1 \cite{iso27033_1}, която разделя мрежата на зони и описва как
се избират мерки за границата между тях. Форматът на сертификата за публичен ключ е дефиниран в
препоръката ITU-T X.509 \cite{itu_x509}, а профилът му за интернет в \cite{rfc5280}. Указанията
на NIST за настройка на TLS \cite{nist80052r2} и за политиката на защитните стени
\cite{nist80041r1} са полезни с това, че формулират проверими твърдения: изискване от вида „да
не се договаря версия под TLS 1.2“ анализаторът на ръкостискането може да провери автоматично.

От класификацията на заплахите в доклада на ENISA \cite{enisa_etl2023} значение тук имат три.
Отказът на услуга, включително разпределеният му вариант. Злоупотребата с инфраструктурата на
публичните ключове, при която проверяващата страна приема сертификат, който не би трябвало да
приеме. И манипулацията на разрешаването на имена: ако името е разрешено до грешен адрес,
проверката на веригата остава единственият механизъм, който може да открие подмяната, защото
удостоверяване на произхода на отговора липсва по устройство \cite{rfc4033}. Класификацията на
OWASP \cite{owasp_top10_2021} е ползвана само като контролен списък какво остава извън обхвата:
проектът работи под приложния слой.

Протоколните спецификации са три. \cite{rfc8446} определя докъде стига пасивният анализ, защото
TLS 1.3 шифрира по-голямата част от ръкостискането. \cite{rfc4987} обяснява защо лавината от
заявки за връзка е евтина за източника и скъпа за получателя, и описва контрамерките заедно със
страничните им ефекти; тези странични ефекти са величината, която проектът измерва.
\cite{rfc9001} показва, че същата ключова обмяна се пренася и в друг транспорт, тоест
механизмите не са свързани с TCP.

\subsection{Приети решения и цената им}
\label{subsec:choices}

Едно решение предопределя останалите: от какво има право да зависи сборката. Обичайният подход
е библиотеките да идват от управител на пакети, при което реализацията е по-кратка, но сборката
изисква подготвена машина и мрежа. Приет е обратният, защото хранилището е публично и трябва да
може да бъде сглобено от непознат без подготовка. Цената е платена на три места: отпадат
библиотеката за прихващане на пакети, външната рамка за тестове и проверката на подписа.

\textbf{Източник на данните.} Анализът отвътре в приложението описва само собствените връзки и е
обвързан с версията на криптографската библиотека. Пасивното прихващане на пакети е общо, но
внася задължителна зависимост. Приетото решение е слоят за разчитане да приема просто изглед към
байтове: общността и възпроизводимостта на входа се запазват, зависимостта отпада. Загубено е
удобството да се чете формат на записан трафик направо, което е отделна задача за разчитане на
формат, а не свойство на анализа.

\textbf{Обхват на собствената реализация.} Криптографските примитиви са зона, в която
собствената реализация е грешка: изборът на алгоритми, устойчивостта срещу атаки по странични
канали и правилността на изчисленията са резултат от години рецензии. Замисълът предвиждаше
примитивите да дойдат от утвърдена библиотека, но това влезе в противоречие с изискването за
сглобяване без подготовка. От трите изхода, а именно задължителна криптографска зависимост,
собствена аритметика с големи числа или непроведена проверка, е приет третият: проверката на
подписа не се извършва и това се обявява при всяка верига. Първите два или отнемат свойството,
заради което хранилището е използваемо, или заменят рецензирана реализация с нерецензирана.
Съществено е, че обявяването е част от изхода: мълчаливото пропускане на проверка прави отчета
по-опасен от липсата на отчет. По същата причина отмяната е сведена до местен списък от серийни
номера, тъй като запитване по \cite{rfc6960} изисква изходяща връзка, каквато системата няма.

\textbf{Стратегия за допускане.} Стойностите без съхранение на състояние \cite{rfc4987}
премахват нуждата от състояние за незавършена връзка, но изискват допълнителен обмен.
Ограничаването по източник е евтино, но е безполезно срещу разпределен източник и вредно, когато
честният и лавинният трафик делят един адрес. Изчислителната задача прехвърля цената върху
клиента, но внася закъснение при първо свързване. Отчитането на таблицата с връзките е по-скоро
измерителен механизъм, но без него останалите три не могат да бъдат оценени. Приети са и
четирите, включвани поотделно, защото въпросът не е кой механизъм е най-добър, а колко струва
всеки, което изисква сравнение при еднакви условия.
